% 来源及取材范围见分题文件；此为整理者转录的正文，不是下载原件。
\begin{theorem}[9.19, L\"uroth]
Let $L=F(X)$ with $X$ transcendental over $F$. Every subfield $E$ of $L$ properly containing $F$ is of the form $E=F(u)$ for some $u\in L$ transcendental over $F$.
\end{theorem}
\paragraph{Gauss's lemma.}
Let $R$ be a unique factorization domain, and let $Q$ be its field of fractions, for example, $R=F[X]$ and $Q=F(X)$. A polynomial $f(T)=\sum a_iT^i$ in $R[T]$ is said to be primitive if its coefficients $a_i$ have no common factor other than units. Every polynomial $f$ in $Q[X]$ can be written $f=c(f)\cdot f_1$ with $c(f)\in Q$ and $f_1$ primitive (write $f=af/a$ with $a$ a common denominator for the coefficients of $f$, and then write $f=(b/a)f_1$ with $b$ the greatest common divisor of the coefficients of $af$). The element $c(f)$ is uniquely determined up to a unit, and $f\in R[X]$ if and only if $c(f)\in R$.

\paragraph{9.20.} If $f,g\in R[T]$ are primitive, so also is $fg$.
\begin{proof}
Let $f=\sum a_iT^i$ and $g=\sum b_iT^i$, and let $p$ be a prime element of $R$. Because $f$ is primitive, there exists a coefficient $a_i$ not divisible by $p$---let $a_{i_1}$ be the first such coefficient. Similarly, let $b_{i_2}$ be the first coefficient of $g$ not divisible by $p$. Then the coefficient of $T^{i_1+i_2}$ in $fg$ is not divisible by $p$. This shows that $fg$ is primitive.
\end{proof}
\paragraph{9.21.} For any $f,g\in R[T]$, $c(fg)=c(f)c(g)$ and $(fg)_1=f_1g_1$.
\begin{proof}
Let $f=c(f)f_1$ and $g=c(g)g_1$ with $f_1$ and $g_1$ primitive. Then $fg=c(f)c(g)f_1g_1$ with $f_1g_1$ primitive, and so $c(fg)=c(f)c(g)$ and $(fg)_1=f_1g_1$.
\end{proof}
\paragraph{9.22.} Let $f$ be a polynomial in $R[T]$. If $f$ factors into the product of two nonconstant polynomials in $Q[T]$, then it factors into the product of two nonconstant polynomials in $R[T]$.
\begin{proof}
Suppose that $f=gh$ in $Q[T]$. Then $f_1=g_1h_1$ in $R[T]$, and $f=c(f)\cdot f_1=(c(f)\cdot g_1)h_1$ is a factorization of $f$ in $R[T]$.
\end{proof}
\paragraph{9.23.} Let $f,g\in R[T]$. If $f$ divides $g$ in $Q[T]$ and $f$ is primitive, then it divides $g$ in $R[T]$.
\begin{proof}
Let $fq=g$ with $q\in Q[T]$. Then $c(q)=c(g)\in R$, and so $q\in R[T]$.
\end{proof}

\paragraph{Proof of L\"uroth's theorem.}
We define the degree $\deg(u)$ of an element $u$ of $F(X)$ to be the larger of the degrees of the numerator and denominator of $u$ when it is expressed in its simplest form.
\begin{lemma}[9.24]
Let $u\in F(X)\setminus F$. Then $u$ is transcendental over $F$, $X$ is algebraic over $F(u)$, and $[F(X):F(u)]=\deg(u)$.
\end{lemma}
\begin{proof}
Let $u(X)=a(X)/b(X)$ with $a(X)$ and $b(X)$ relatively prime polynomials. Then $a(T)-b(T)u$ is a polynomial in $F(u)[T]$ having $X$ as a root, and so $X$ is algebraic over $F(u)$. It follows that $u$ is transcendental over $F$ (else $X$ would be algebraic over $F$; 1.31).

The polynomial $a(T)-b(T)Z\in F[Z,T]$ is clearly irreducible. As $u$ is transcendental over $F$,
\[
F[Z,T]\simeq F[u,T],\qquad Z\leftrightarrow u,\quad T\leftrightarrow T,
\]
and so $a(T)-b(T)u$ is irreducible in $F[u,T]$, and hence also in $F(u)[T]$ by Gauss's lemma (9.22). It has $X$ as a root, and so, up to a constant, it is the minimal polynomial of $X$ over $F(u),$ and its degree is $\deg(u)$, which proves the lemma.
\end{proof}

\begin{proof}[Proof of Theorem 9.19]
For any $u\in E\setminus F$,
\[
[F(X):E]\leq[F(X):F(u)]=\deg(u),
\]
and so $X$ is algebraic over $E$. Let
\[
f(T)=T^n+a_1T^{n-1}+\cdots+a_n,\qquad a_i\in E,\quad n=[F(X):E],
\]
be its minimal polynomial. As $X$ is transcendental over $F$, some $a_j\notin F$, and we'll show that $E=F(a_j)$ for such an $a_j$.

Let $d(X)\in F[X]$ be a polynomial of least degree such that $d(X)a_i(X)\in F[X]$ for all $i$, and let
\[
f_1(X,T)=df(T)=dT^n+da_1T^{n-1}+\cdots+da_n\in F[X,T].
\]
Then $f_1$ is primitive as a polynomial in $T$, i.e., $\gcd(d,da_1,\ldots,da_n)=1$ in $F[X]$. The degree $m$ of $f_1$ in $X$ is the largest degree of one of the polynomials $da_1,da_2,\ldots$, say, $m=\deg(da_i)$.

Write $a_i=b/c$ with $b,c$ relatively prime polynomials in $F[X]$. Now $b(T)-c(T)a_i(X)$ is a polynomial in $E[T]$ having $X$ as a root, and so it is divisible by $f$, say
\[
f(T)\cdot q(T)=b(T)-c(T)\cdot a_i(X),\qquad q(T)\in E[T].
\]
On multiplying through by $c(X)$, we find that
\[
c(X)\cdot f(T)\cdot q(T)=c(X)\cdot b(T)-c(T)\cdot b(X).
\]
As $f_1$ differs from $f$ by a nonzero element of $F(X)$, the equation shows that $f_1$ divides $c(X)\cdot b(T)-c(T)\cdot b(X)$ in $F(X)[T]$, but $f_1$ is primitive in $F[X][T]$, and so it divides the polynomial in $F[X][T]=F[X,T]$ (by 9.23), i.e., there exists a polynomial $h\in F[X,T]$ such that
\[
f_1(X,T)\cdot h(X,T)=c(X)\cdot b(T)-c(T)\cdot b(X).\tag{14}
\]
In (14), the polynomial $c(X)\cdot b(T)-c(T)\cdot b(X)$ has degree at most $m$ in $X$, and $m$ is the degree of $f_1(X,T)$ in $X$. Therefore, $c(X)\cdot b(T)-c(T)\cdot b(X)$ has degree exactly $m$ in $X$, and $h(X,T)$ has degree $0$ in $X$, i.e., $h\in F[T]$. It now follows from (14) that $c(X)\cdot b(T)-c(T)\cdot b(X)$ is not divisible by a nonconstant polynomial in $F[X]$.

The polynomial $c(X)\cdot b(T)-c(T)\cdot b(X)$ is symmetric in $X$ and $T$, i.e., it is unchanged when they are swapped. Therefore, it has degree $m$ in $T$ and it is not divisible by a nonconstant polynomial in $F[T]$. It now follows from (14) that $h$ is not divisible by a nonconstant polynomial in $F[T]$, and so it lies in $F^\times$. We conclude that $f_1(X,T)$ is a constant multiple of $c(X)\cdot b(T)-c(T)\cdot b(X)$.

On comparing degrees in $T$ in (14), we see that $n=m$. Thus
\begin{align*}
[F(X):F(a_i)]&\overset{9.24}{=}\deg(a_i)\leq\deg(da_i)=m=n\,;\\
n&=[F(X):E]\leq[F(X):F(a_i)].
\end{align*}
Hence, equality holds throughout, and so $E=F[a_i]$.
Finally, if $a_j\notin F$, then
\begin{align*}
[F(X):E]&\leq[F(X):F(a_j)]\overset{9.24}{=}\deg(a_j)\\
&\leq\deg(da_j)\leq\deg(da_i)=m=[F(X):E],
\end{align*}
and so $E=F(a_j)$ as claimed.
\end{proof}
